% Преамбула книги о прошивке LiteBee. Сборка: latexmk -xelatex kniga.tex
\usepackage{fontspec}
\usepackage{polyglossia}
\setmainlanguage{russian}
\setotherlanguage{english}
\setmainfont{Charis SIL}[Scale=0.93]
\setsansfont{Fira Sans}
\setmonofont{JetBrains Mono}[Scale=0.84, RawFeature={-calt;-liga}]
\newfontfamily\cyrillicfonttt{JetBrains Mono}[Scale=0.84, RawFeature={-calt;-liga}]

\usepackage{xcolor}
\definecolor{accent}{HTML}{0B7285}
\definecolor{accentdark}{HTML}{095C6B}
\definecolor{accentsoft}{HTML}{E3F1F4}
\definecolor{ink}{HTML}{1F2933}
\definecolor{inkgray}{HTML}{5B6770}
\definecolor{rule}{HTML}{C9D2D6}
\definecolor{paper}{HTML}{F5F7F8}
\definecolor{warnfill}{HTML}{FBF3E4}
\definecolor{warnedge}{HTML}{B7791F}
\definecolor{stopfill}{HTML}{FAECEC}
\definecolor{stopedge}{HTML}{A33A3A}
\definecolor{ledred}{HTML}{D64545}
\definecolor{ledgreen}{HTML}{3C9D5A}
\definecolor{ledblue}{HTML}{3465C4}
\definecolor{ledyellow}{HTML}{D9B22B}
\definecolor{ledpurple}{HTML}{8A5CC4}
\definecolor{ledcyan}{HTML}{2AA0B8}

\usepackage{geometry}
\geometry{a4paper, top=26mm, bottom=28mm, inner=24mm, outer=24mm, headsep=9mm}
\KOMAoptions{parskip=half-}
\linespread{1.06}

\usepackage{microtype}
\usepackage{graphicx}
\usepackage{booktabs}
\usepackage{longtable}
\usepackage{tabularx}
\usepackage{array}
\usepackage{enumitem}
\setlist{itemsep=2pt, topsep=4pt, parsep=0pt}
\setlist[enumerate,1]{leftmargin=1.9em}
\setlist[itemize,1]{leftmargin=1.3em}
\usepackage{caption}
\captionsetup{font={small,color=inkgray}, labelfont={bf,color=ink}, labelsep=period, justification=raggedright, singlelinecheck=false, skip=6pt}
\usepackage{fvextra}
\usepackage[most]{tcolorbox}
\usepackage{tikz}
\usetikzlibrary{arrows.meta, calc, positioning, fit, shapes.geometric, backgrounds, decorations.pathreplacing}
\usepackage{float}
\usepackage{needspace}
\floatplacement{figure}{tbp}
\raggedbottom

% ---------- заголовки
\addtokomafont{disposition}{\sffamily\color{ink}}
\KOMAoptions{headings=small}
\renewcommand*{\chaptername}{Глава}
\setkomafont{chapter}{\sffamily\bfseries\fontsize{24}{28}\selectfont\color{accent}}
\setkomafont{chapterprefix}{\sffamily\bfseries\normalsize\color{inkgray}}
\RedeclareSectionCommand[beforeskip=0pt, afterskip=1.6\baselineskip, innerskip=4pt]{chapter}
\renewcommand*{\chapterformat}{\chapapp\nobreakspace\thechapter}
\renewcommand*{\chapterlinesformat}[3]{%
  \parbox{\textwidth}{\ifstr{#2}{}{}{{\usekomafont{chapterprefix}#2}\par\vspace{2pt}}{#3}\par\vspace{6pt}%
  \textcolor{accent}{\rule{3.2cm}{1.6pt}}}}
\setkomafont{section}{\sffamily\bfseries\Large\color{accentdark}}
\setkomafont{subsection}{\sffamily\bfseries\large\color{ink}}
\setkomafont{subsubsection}{\sffamily\bfseries\normalsize\color{ink}}
\RedeclareSectionCommand[beforeskip=1.5\baselineskip plus 4pt, afterskip=.5\baselineskip]{section}
\RedeclareSectionCommand[beforeskip=1.1\baselineskip plus 3pt, afterskip=.3\baselineskip]{subsection}
\setkomafont{paragraph}{\sffamily\bfseries\normalsize\color{ink}}
\RedeclareSectionCommand[beforeskip=.7\baselineskip, afterskip=-.5em]{paragraph}

\usepackage[automark,headsepline,plainheadsepline]{scrlayer-scrpage}
\pagestyle{scrheadings}
\clearpairofpagestyles
\ihead{\headmark}
\ohead{\pagemark}
\setkomafont{pagehead}{\sffamily\small\color{inkgray}}
\setkomafont{pagenumber}{\sffamily\small\bfseries\color{ink}}
\setkomafont{headsepline}{\color{rule}}
\renewcommand*{\chaptermarkformat}{\chapapp\nobreakspace\thechapter.\enskip}

\usepackage[unicode, hidelinks, pdftitle={Как устроены дроны LiteBee}, pdfauthor={Litebee}]{hyperref}
\hypersetup{colorlinks=true, linkcolor=accent, urlcolor=accent, citecolor=accent}

% ---------- листинги
\fvset{fontsize=\small, breaklines=true, breakanywhere=false, breaksymbolleft={}, breakindent=0pt,
  breakautoindent=false, numbers=left, numbersep=7pt, xleftmargin=16pt, tabsize=4}
\renewcommand{\theFancyVerbLine}{\textcolor{inkgray}{\scriptsize\arabic{FancyVerbLine}}}

\tcbset{
  codebox/.style={enhanced, colback=paper, colframe=paper, boxrule=0pt, sharp corners, frame hidden,
    borderline west={2.5pt}{0pt}{accent}, left=9pt, right=6pt, top=3pt, bottom=3pt,
    before skip=2pt, after skip=10pt},
  note/.style={enhanced, breakable, colback=accentsoft, colframe=accentsoft, frame hidden, sharp corners,
    borderline west={3pt}{0pt}{accent}, left=9pt, right=8pt, top=5pt, bottom=5pt, before skip=8pt, after skip=10pt},
  warn/.style={enhanced, breakable, colback=warnfill, colframe=warnfill, frame hidden, sharp corners,
    borderline west={3pt}{0pt}{warnedge}, left=9pt, right=8pt, top=5pt, bottom=5pt, before skip=8pt, after skip=10pt},
  stop/.style={enhanced, breakable, colback=stopfill, colframe=stopfill, frame hidden, sharp corners,
    borderline west={3pt}{0pt}{stopedge}, left=9pt, right=8pt, top=5pt, bottom=5pt, before skip=8pt, after skip=10pt},
  try/.style={enhanced, breakable, colback=white, colframe=accent, boxrule=0.9pt, arc=3pt,
    left=10pt, right=10pt, top=6pt, bottom=6pt, before skip=12pt, after skip=10pt,
    fonttitle=\sffamily\bfseries\normalsize, coltitle=white, colbacktitle=accent, toptitle=3pt, bottomtitle=3pt},
}

% подпись файла над листингом + сам листинг; #1 — настройки VerbatimInput, #2 — путь
\newcommand{\filelabel}[1]{\par\needspace{5\baselineskip}\noindent{\sffamily\scriptsize\color{inkgray}\texttt{\detokenize{#1}}}\par\nopagebreak\vspace{1pt}}
\newcommand{\listing}[2][]{%
  \filelabel{#2}%
  \begin{tcolorbox}[codebox]\VerbatimInput[#1]{#2}\end{tcolorbox}}
\newcommand{\solfile}[2][]{\listing[#1]{solutions/#2}}
% длинный листинг без рамки, для ответов
\newcommand{\answerfile}[2][]{%
  \filelabel{solutions/#2}%
  \begin{tcolorbox}[codebox]\VerbatimInput[fontsize=\footnotesize,#1]{solutions/#2}\end{tcolorbox}}

% ---------- выноски
\newenvironment{zametka}[1][Запомните]{\begin{tcolorbox}[note]{\sffamily\bfseries\small\color{accentdark}#1}\par\smallskip}{\end{tcolorbox}}
\newenvironment{vnimanie}[1][Осторожно]{\begin{tcolorbox}[warn]{\sffamily\bfseries\small\color{warnedge}#1}\par\smallskip}{\end{tcolorbox}}
\newenvironment{opasno}[1][Безопасность]{\begin{tcolorbox}[stop]{\sffamily\bfseries\small\color{stopedge}#1}\par\smallskip}{\end{tcolorbox}}

% ---------- короткие обозначения в тексте
\newcommand{\kod}[1]{\texttt{\detokenize{#1}}}
\newcommand{\knopka}[1]{\textsf{«#1»}}
\newcommand{\klavisha}[1]{\tcbox[on line, colback=paper, colframe=rule, boxrule=0.4pt, arc=2pt, boxsep=1pt, left=2pt, right=2pt, top=0pt, bottom=0pt]{\sffamily\small #1}}
\usepackage{xspace}
\newcommand{\cm}{\,см\xspace}
\newcommand{\sek}{\,с\xspace}
\addto\captionsrussian{\renewcommand*{\chaptername}{Глава}}
\newcommand{\status}[1]{\textsf{\small #1}}

% снимок окна: обрезанный фрагмент на всю ширину текста
\newcommand{\snimok}[3]{%
  \begin{figure}[tbp]
    \centering
    \fbox{\includegraphics[width=\dimexpr\textwidth-2\fboxsep-2\fboxrule\relax]{figures/#1.png}}
    \caption{#2}\label{#3}
  \end{figure}}

\setlength{\fboxsep}{0pt}
\setlength{\fboxrule}{0.4pt}
\hfuzz=1.5pt
\emergencystretch=3em
\tolerance=1500
\usepackage{url}\urlstyle{tt}
\DeclareTOCStyleEntry[indent=0pt, numwidth=2.4em]{tocline}{chapter}
\DeclareTOCStyleEntry[indent=2.4em, numwidth=3em]{tocline}{section}
\setcounter{tocdepth}{1}
\setcounter{secnumdepth}{1}
\renewcommand{\arraystretch}{1.25}

% ---------- метки достоверности
\newcommand{\metka}[2]{\tcbox[on line, colback=#1!14, colframe=#1!70, boxrule=0.4pt, arc=2pt, boxsep=0.5pt, left=2.5pt, right=2.5pt, top=0pt, bottom=0pt]{\sffamily\scriptsize\bfseries\color{#1!80!black}#2}}
\newcommand{\pr}{\metka{ledgreen}{ПРОЧИТАНО}}
\newcommand{\vyv}{\metka{ledyellow}{ВЫВОД}}
\newcommand{\nei}{\metka{ledred}{НЕИЗВЕСТНО}}
\newcommand{\adr}[1]{\texttt{\detokenize{#1}}}
\newcolumntype{Y}{>{\raggedright\arraybackslash}X}
\newcommand{\hdr}[1]{\textbf{\sffamily\small #1}}
\renewcommand{\arraystretch}{1.18}
